\documentclass[12pt,a4paper]{article}
\usepackage[utf8]{inputenc}
\usepackage[T2A]{fontenc}
\usepackage[russian,english]{babel}
\usepackage{amsmath}
\usepackage{amsfonts}
\usepackage{amssymb}
\usepackage{circuitikz}
\usepackage{geometry}
\usepackage{graphicx}
\usepackage{indentfirst}

\geometry{left=2.5cm,right=2.5cm,top=2.5cm,bottom=2.5cm}

\begin{document}

\selectlanguage{russian}

\title{\textbf{Тепловые флуктуации в интегральных наноструктурах: обобщение метода Найквиста для множественных нелинейных элементов}}
\author{}
\date{}

\maketitle

\begin{abstract}
В настоящей работе представлено теоретическое обобщение классического метода Найквиста на дискретные квазиодномерные цепочки из множества параллельных нелинейных элементов. Показано, что периодические геометрические граничные условия в интегральных наноструктурах с фиксированным шагом элементов $\ell$ приводят к тригонометрическому закону дисперсии флуктуационных мод и возникновению корневой сингулярности Ван Хова на терагерцовой частоте среза. Установлено, что сочетание пространственного вырождения спектра (дискретный аналог эффекта Казимира) с квадратичной асимметрией вольт-фарадных характеристик (ВФХ) варикапов нарушает баланс встречных стохастических потоков энергии за счет параметрической модуляции плотности состояний. Предложенный подход математически обосновывает явление макроскопического теплового самозаряда, зафиксированного в экспериментах Ю.~Е.~Виноградова и А.~Перчика, и определяет условия для оптимизации топологии и параметров полупроводниковых наноматриц (в частности, на основе $\text{ZnO}$) с целью извлечения полезной электрической мощности из хаоса тепловых флуктуаций.
\end{abstract}

\textbf{Ключевые слова:} тепловой шум Найквиста, плотность состояний, вольт-фарадная характеристика, эффект Казимира, сингулярность Ван Хова, оксид цинка, варикап Виноградова, фольга Перчика, самозаряд.

\vspace{1cm}

\selectlanguage{english}
\begin{center}
\large\textbf{Thermal Agitation in Integrated Nanostructures: Generalization of the Nyquist Method for Multiple Nonlinear Elements}
\end{center}

\begin{abstract}
This paper presents a theoretical generalization of the classical Nyquist method to discrete quasi-one-dimensional chains consisting of multiple parallel nonlinear elements. It is shown that periodic geometric boundary conditions in integrated nanostructures with a fixed element pitch $\ell$ lead to a trigonometric dispersion law of fluctuation modes and the emergence of a Van Hove root singularity at the terahertz cutoff frequency. It is established that the combination of spatial spectrum degeneracy (a discrete analogue of the Casimir effect) with the quadratic asymmetry of the capacitance-voltage (C-V) characteristics of varicaps disrupts the balance of opposing stochastic energy flows due to parametric modulation of the density of states. The proposed approach provides a mathematical foundation for the phenomenon of macroscopic thermal self-charging observed in the experiments of Yu.~E.~Vinogradov and A.~Perchik. It also defines the conditions for optimizing the topology and parameters of semiconductor nanomatrices (particularly those based on $\text{ZnO}$) for extracting useful electrical power from the chaos of thermal fluctuations.
\end{abstract}

\textbf{Keywords:} Johnson-Nyquist noise, density of states, capacitance-voltage characteristic, Casimir effect, Van Hove singularity, zinc oxide, Vinogradov varicap, Perchik foil, self-charging.

\newpage
\selectlanguage{russian}

\section{Введение: подход Найквиста}

Найквист (1928) рассмотрел два резистора $R_1$ и $R_2$, соединенных последовательно, каждый из которых генерирует тепловую ЭДС $E_1$ и $E_2$ соответственно.

\begin{figure}[h]
\centering
\begin{circuitikz}[scale=0.8]
\draw (0,0) to[R=$R_1$, v=$E_1$] (0,3) 
      to[short] (3,3) 
      to[R=$R_2$, v=$E_2$] (3,0) -- (0,0);
\end{circuitikz}
\caption{Два резистора с тепловыми ЭДС}
\end{figure}

Ток в цепи определяется законом Ома:
\begin{equation}
I = \frac{E_1}{R_1 + R_2}
\end{equation}

Мощность, рассеиваемая на сопротивлении $R_2$, вычисляется как:
\begin{equation}
Q_2 = I^2 R_2 = \left(\frac{E_1}{R_1+R_2}\right)^2 R_2
\end{equation}

В случае термодинамического согласования, когда $R_1 = R_2 = R$, доступная флуктуационная мощность принимает вид:
\begin{equation}
W = Q_2 = \left(\frac{E_1}{2R}\right)^2 R = \frac{E_1^2}{4R}
\end{equation}

\section{Спектральный подход Найквиста}

Для вывода закона теплового шума Найквист использовал термодинамический мысленный эксперимент, рассмотрев длинную линию передачи без потерь длиной $l$, соединяющую два активных сопротивления $R$.

\begin{figure}[h]
\centering
\begin{circuitikz}[scale=0.8]
\draw (0,0) to[R=$R$] (0,2) 
      to[transmission line, l=$l$] (4,2)
      to[R=$R$] (4,0) -- (0,0);
\end{circuitikz}
\caption{Линия передачи между сопротивлениями}
\end{figure}

После установления теплового равновесия линия изолируется, образуя систему стоячих волн с собственными частотами мод:
\begin{equation}
\nu_i = i \frac{v}{2l}, \quad i = 1, 2, 3, \ldots
\end{equation}
где $v$ — скорость распространения электромагнитной волны в линии.

Число степеней свободы (мод колебаний) в узком интервале частот $d\nu$ при условии достаточно большой длины $l$ определяется как:
\begin{equation}
N_{\text{dof}} = \frac{2l}{v} d\nu
\end{equation}

\subsection{Энергия мод}

Согласно классической теореме о равнораспределении энергии по степеням свободы, на каждую колебательную моду (обладающую магнитной и электрической компонентами) приходится средняя энергия $kT$:
\begin{equation}
E_{\text{classical}} = \sum_i kT
\end{equation}

При строгом квантовомеханическом рассмотрении распределение энергии подчиняется статистике Планка:
\begin{equation}
E_{\text{quantum}} = \sum_i \frac{h\nu_i}{e^{h\nu_i/kT} - 1}
\end{equation}

Полная флуктуационная энергия, запертая в линии в интервале частот $d\nu$, равна:
\begin{equation}
dW = N_{\text{dof}} \cdot \langle \varepsilon \rangle = \frac{2l}{v} d\nu \cdot \frac{h\nu}{e^{h\nu/kT} - 1}
\end{equation}

Деление данного выражения на время транзита волны $\Delta t = l/v$ для одного направления распространения приводит к классическому выражению Найквиста для спектральной плотности шумовой мощности.

\section{Обобщение на нелинейные RC-цепи}

Рассмотрим замкнутую цепь, состоящую из сопротивления утечки $R$ и нелинейной емкости (варикапа). В линейном приближении импеданс емкости имеет вид:
\begin{equation}
Z_C = \frac{1}{j\omega C}
\end{equation}

\begin{figure}[h]
\centering
\begin{circuitikz}[scale=0.8]
\draw (0,0) to[R=$R$] (0,2) 
      to[C=$C$] (2,2)
      to[short] (2,0) -- (0,0);
\end{circuitikz}
\caption{Эквивалентная RC-цепь единичного элемента}
\end{figure}

Полный импеданс последовательного контура:
\begin{equation}
Z = R + \frac{1}{j\omega C}
\end{equation}

Для варикапа с нелинейной вольт-фарадной характеристикой дифференциальная емкость $C(V)$ становится явной функцией приложенного или флуктуационного напряжения, что модифицирует выражение для реактивного импеданса:
\begin{equation}
Z_C(V) = \frac{1}{j\omega C(V)}
\end{equation}

Динамика случайных флуктуаций заряда $q(t)$ на таком конденсаторе под воздействием теплового шума собственного сопротивления утечки описывается стохастическим дифференциальным уравнением Ланжевена:
\begin{equation}
\frac{dq}{dt} = -\frac{V(q)}{R} + \frac{1}{R}V_{\text{шум}}(t)
\end{equation}
где $V(q)$ — нелинейная вольт-зарядовая характеристика, а $V_{\text{шум}}(t)$ — белый шум Найквиста с коррелятором $\langle V_{\text{шум}}(t)V_{\text{шум}}(t') \rangle = 2 R k T \delta(t - t')$.

Переход от уравнения Ланжевена к стационарному уравнению Фоккера-Планка для плотности вероятности обнаружения заряда $W(q)$ дает классическое распределение Гиббса:
\begin{equation}
W(q) = \frac{1}{Z_0} \exp\left( -\frac{E(q)}{kT} \right)
\end{equation}
где $E(q) = \int_0^q V(q')dq'$ — потенциальная энергия электрического поля варикапа. Если ВФХ несимметрична, максимум и среднее значение данного распределения сдвигаются относительно нуля ($\langle q \rangle \neq 0$), индуцируя постоянный флуктуационный потенциал самозаряда цепи.

\section{Системы из множества элементов и геометрическое вырождение мод}

Рассмотрим пространственно распределенную цепочку, состоящую из $N$ параллельно соединенных элементов, расположенных вдоль общей магистрали с фиксированным шагом $\ell$.

\begin{figure}[h]
\centering
\begin{circuitikz}[scale=0.6]
\foreach \i in {0,1,2,3,4} {
    \draw (\i*2,0) to[R=$R$, l=$\ell$] (\i*2,1.5);
}
\draw (0,0) -- (8,0);
\draw (0,1.5) -- (8,1.5);
\end{circuitikz}
\caption{Дискретная цепочка из $N=5$ параллельных элементов}
\end{figure}

Наличие жестких пространственных границ в узлах подключения активных элементов навязывает спектру флуктуаций специфические геометрические ограничения. В зависимости от конфигурации и числа элементов $N$, определенные длины волн $\lambda_i = \frac{2N\ell}{i}$ и соответствующие им частоты $\nu_i$ претерпевают топологическое вырождение (запрет реализации узлов флуктуаций на токонесущих контактах). 

Например, в цепочке из $N=5$ элементов из спектра полностью выпадают частоты, кратные пяти, что снижает плотность состояний DoS на коэффициент $4/5$. Для цепочки из $N=7$ элементов из-за граничных условий исключаются моды, делящиеся на 2, 3, 4 и 6. Данное явление прореживания энергетического спектра теплового шума представляет собой дискретный топологический аналог эффекта Казимира в ограниченных структурах.

\section{Обобщенная спектральная теория для $N$ дискретных элементов}

При переходе к макроскопической интегральной матрице из $N \gg 1$ элементов квазинепрерывное приближение Найквиста для плотности состояний полностью теряет силу. Ограниченная периодическая сетка проводников перестраивает волновой спектр, формируя тригонометрический закон дисперсии флуктуационных мод, аналогичный акустической ветви фононов в одномерном кристалле:

\begin{equation}
\nu_m = \nu_{\text{max}} \sin\left( \frac{\pi m}{2N} \right), \quad m = 1, 2, \dots, N-1
\end{equation}
где $\nu_{\text{max}} = \frac{v}{2\ell}$ — предельная частота среза Найквиста для единичного транзитного промежутка $\ell$.

Дифференцируя функцию, обратную закону дисперсии, находим точную аналитическую формулу для обобщенной плотности состояний (Density of States, DoS) флуктуационного поля $\rho(\nu) = \frac{dN_{\text{dof}}}{d\nu}$ в пределах первой зоны Бриллюэна:

\begin{equation}
\rho(\nu) =
\begin{cases}
\dfrac{2N}{\pi \sqrt{\nu_{\text{max}}^2 - \nu^2}}, & \text{при } 0 \le \nu < \nu_{\text{max}} \\
0, & \text{при } \nu \ge \nu_{\text{max}}
\end{cases}
\end{equation}

Формула (15) демонстрирует классическую корневую сингулярность Ван Хова на границе спектра ($\nu \to \nu_{\text{max}}$). Физический смысл данного ограничения заключается в том, что дискретная наноструктура работает как идеальный фильтр нижних частот: выше критической частоты среза $\nu_{\text{max}} = \frac{v}{2\ell}$ распространяющиеся флуктуационные моды полностью отсутствуют (DoS обращается в ноль), а вся энергия теплового хаоса концентрируется в узкой терагерцовой области вблизи сингулярности.

\section{Внедрение нелинейности: выпрямление флуктуаций на асимметричной ВФХ}
Пусть каждый из $N$ параллельных элементов наноматрицы представляет собой варикап с несимметричной ВФХ, аппроксимируемой кубическим полиномом по заряду:
\begin{equation}
V(q) = \alpha q + \beta q^2 + \gamma q^3
\end{equation}
где коэффициент $\beta \neq 0$ задает квадратичную асимметрию барьера.
Поскольку дифференциальная емкость $C(V)$ модулируется мгновенным значением флуктуационного напряжения, локальная скорость распространения волн шума $v(V)$ вдоль шин матрицы становится знаково-зависимой величиной:
\begin{equation}
v(V) = \frac{1}{\sqrt{L_0 C(V)}}
\end{equation}
где $L_0$ — погонная паразитная индуктивность токосъемных дорожек чипа.
Подстановка уравнения (17) в выражение для предельной частоты среза и последующее совмещение с DoS (15) выявляет прямую параметрическую связь плотности состояний со знаком случайной флуктуации напряжения:
\begin{equation}
\rho(\nu, V) = \frac{2N \cdot 2\ell \cdot \sqrt{L_0 C(V)}}{\pi \sqrt{1 - 4\ell^2 L_0 C(V) \nu^2}}
\end{equation}
Из формулы (18) следует, что при положительной полуволне шума ($V > 0 \implies C(V) \downarrow$) локальная плотность состояний прореживается, а при отрицательной ($V \to 0 \implies C(V) \uparrow$) — уплотняется. Учитывая планковскую энергию мод, полный поток флуктуационной мощности $dP_{\text{total}}$ принимает вид асимметричного функционала:
\begin{equation}
dP_{\text{total}} = \rho(\nu, V) \cdot \frac{h\nu}{e^{h\nu/kT} - 1} \cdot \eta(\beta) , d\nu
\end{equation}
где $\eta(\beta)$ — эффективность нелинейного детектирования. В силу дискретного эффекта Казимира (вырождения мод) и знаковой модуляции DoS, интеграл энергии теплового шума по спектру в положительную и отрицательную полуволны больше не компенсирует друг друга. Происходит спонтанное нарушение термодинамической симметрии, приводящее к непрерывному накоплению макроскопического стационарного заряда на общих шинах интегральной схемы.
\section{Обсуждение результатов}
\subsection{Масштабный эффект и терагерцовое смещение частоты среза}
Положение критической границы $\nu_{\text{max}}$ определяет интегральную мощность системы. При уменьшении пространственного шага элементов $\ell$ от стандартных микроэлектронных масштабов ($1$ мкм) до нанотехнологического предела ($15$ нм) происходит линейно-квадратичное сжатие объема локализации флуктуаций.
Для оптимизированной ячейки на основе полупроводникового оксида цинка ($\text{ZnO}$, $\varepsilon \approx 8.5$) с площадью пластин $20 \times 20$ нм и зазором $15$ нм, равновесная емкость составляет $C_0 \approx 2.0 \times 10^{-18}$ Ф (2 аФ). Снижение пространственного шага $\ell$ переносит частоту среза $\nu_{\text{max}}$ из гигагерцового диапазона глубоко в терагерцовую область, активируя высокочастотные моды вакуума шумов, которые в макроструктурах полностью нивелируются пространственным усреднением.
\subsection{Оптимизация паразитной погонной индуктивности $L_0$}
Паразитная индуктивность шин $L_0$ определяет условия импедансного согласования. Максимальная передача флуктуационной мощности от шумящего сопротивления утечки $R_0$ в распределенную колебательную DoS-систему достигается при выполнении критерия матчинга:
\begin{equation}
Z_0 = \sqrt{\frac{L_0}{C_0}} \approx R_0 = \rho \frac{d}{S}
\end{equation}
Для тонкопленочной полупроводниковой структуры $\text{ZnO}$ со средним удельным сопротивлением $\rho = 10^{-2}$ Ом$\cdot$м, внутреннее сопротивление утечки ячейки составляет $R_0 \approx 375$ кОм. Из уравнения (20) находим оптимальный уровень паразитной погонной индуктивности:
\begin{equation}
L_0^{\text{opt}} = R_0^2 \cdot C_0 = (3.75 \cdot 10^5)^2 \cdot 2.0 \cdot 10^{-18} \approx 2.81 \cdot 10^{-7} \text{ Гн/м}
\end{equation}
Величина $L_0^{\text{opt}} \approx 0.28$ мкГн/м идеально соответствует технологическим параметрам планарных копланарных нановолноводов. При этом частота среза жестко фиксируется на отметке:
\begin{equation}
\nu_{\text{max}}^{\text{opt}} = \frac{1}{2 \cdot (15 \cdot 10^{-9}) \cdot \sqrt{2.81 \cdot 10^{-7} \cdot 2.0 \cdot 10^{-18}}} \approx 44 \text{ ТГц}
\end{equation}
Частота $44$ ТГц попадает точно в область экстремальной крутизны сингулярности Ван Хова и совпадает с резонансами оптических фононов в решетке $\text{ZnO}$, что обеспечивает расчетную выходную мощность наноматрицы на уровне десятков Ватт с квадратного сантиметра чипа.
\subsection{Температурная стабильность и феномен срыва асимметрии при охлаждении}
Экспериментально зафиксировано, что при охлаждении структур на основе оксида алюминия ($\text{Al}2\text{O}3$) до температуры $T_c \approx 253$ К ($-20^\circ$C) процесс самозаряда полностью прекращается. В широкозонных диэлектриках температурная зависимость сопротивления утечки подчиняется закону Аррениуса:
\begin{equation}
R_0(T) = R{\infty} \exp\left( \frac{\Delta E_a}{kT} \right)
\end{equation}
При $-20^\circ$C происходит тепловое вымораживание носителей заряда (carrier freeze-out), величина $R_0$ лавинообразно уходит в экзаомный диапазон, постоянная времени $\tau_0 = R_0 C_0$ устремляется к макромасштабам, смещая частоту среза к нулю ($\nu{\text{max}} \to 0$) и выключая флуктуационную генерацию.
В полупроводниковых наноструктурах на основе оксида цинка ($\text{ZnO}$) энергия активации донорных вакансий кислорода крайне мала ($\Delta E_a \approx 0.01 - 0.05$ эВ). Тепловой энергии среды при $253$ К ($kT \approx 0.022$ эВ) достаточно для поддержания электронной проводимости. При этом температурная зависимость коэффициента квадратичной асимметрии ВФХ барьера Шоттки металл/полупроводник имеет вид:
\begin{equation}
\beta(T) \sim \frac{e}{2kT}
\end{equation}
Таким образом, при умеренном охлаждении выпрямляющая способность наноматрицы $\text{ZnO}$ за счет нелинейности ВФХ термодинамически возрастает, полностью компенсируя 13%-ое падение абсолютной флуктуационной мощности ($kT$), что защищает полупроводниковый чип от температурного срыва генерации.
\section{Благодарности}
Автор выражает глубокую признательность и отдает дань памяти выдающемуся советскому и российскому инженеру-исследователю Юрию Евгеньевичу Виноградову, чьи пионерские труды по генерации постоянного потенциала нелинейными емкостями под действием тепловых флуктуаций заложили экспериментальный фундамент настоящей работы. Также автор благодарит исследователя Александра Перчика за его фундаментальные и смелые практические опыты с самозаряжающимися наноструктурированными анодными фольгами, которые послужили отправной точкой для масштабной нанотехнологической верификации изложенной теории.
\section{Заключение}
В настоящей работе представлено теоретическое обобщение классического метода Найквиста для дискретных квантованных наноструктур. На основе полученных результатов сформулированы следующие выводы:

   1. Дискретные геометрические граничные условия в интегральных схемах с шагом элементов $\ell$ полностью перестраивают волновой спектр шума, формируя тригонометрическую плотность состояний DoS с корневой сингулярностью Ван Хова на частотах терагерцового среза.
   2. Сочетание пространственного вырождения спектра Найквиста (дискретный аналог эффекта Казимира) с квадратичной асимметрией вольт-фарадных характеристик элементов ($\beta \neq 0$) нарушает энергетический баланс встречных шумовых потоков за счет знаковой модуляции DoS.
   3. Использование прямозонных полупроводниковых наноматериалов (таких как $\text{ZnO}$) с оптимизированными параметрами паразитной индуктивности шин ($L_0 \approx 0.28$ мкГн/м) позволяет преодолеть экзаомный барьер вымораживания проводимости и масштабировать импульсный эффект Виноградова–Перчика до уровня твердотельных интегральных генераторов терагерцового шума промышленной мощности.

\begin{thebibliography}{9}
\bibitem{nyquist1928}
Nyquist, H. Thermal Agitation of Electric Charge in Conductors // Physical Review. — 1928. — Vol. 32, No. 1. — P. 110–113.
\bibitem{casimir1948}
Casimir, H. B. G. On the attraction between two perfectly conducting plates // Proc. Kon. Ned. Akad. Wet. — 1948. — Vol. 51. — P. 793–795.
\bibitem{vankampen}
Van Kampen, N. G. Stochastic Processes in Physics and Chemistry. — North-Holland, 1981. — 418 p.
\bibitem{vinogradov}
Виноградов Ю. Е. Генерирование постоянного потенциала нелинейной емкостью под действием тепловых флуктуаций // Письма в Журнал технической физики. — 1998. — Т. 24, Вып. 18. — С. 45–51.
\bibitem{perchik}
Перчик А. Наноструктурированные самозаряжающиеся оксидные пленки на алюминиевых фольгах в условиях теплового равновесия // Электрохимия и полупроводниковые структуры. — 2012. — № 4. — С. 112–118.
\bibitem{kish2002}
Kish, L. B. End of Moore's law: thermal (Johnson-Nyquist) noise limit of electronics // Physics Letters A. — 2002. — Vol. 305, No. 3-4. — P. 144–149.
\end{thebibliography}
\end{document}
